\documentclass[11pt,a4paper]{article}

\usepackage{fontspec}
\usepackage{polyglossia}
\setdefaultlanguage{english}

\usepackage[a4paper,margin=1.8cm]{geometry}
\usepackage{amsmath,amssymb}
\usepackage{booktabs}
\usepackage{longtable}
\usepackage{tabularx}
\usepackage{array}
\usepackage{graphicx}
\usepackage{float}
\usepackage{tikz}
\usetikzlibrary{arrows.meta, positioning, shapes}
\usepackage{sectsty}
\usepackage{pgfplots}
\usepackage{pgfplotstable}
\usepackage{xcolor}
\usepackage{microtype}
\usepackage{enumitem}
\usepackage{hyperref}
\usepackage{fancyhdr}
\usepackage{titlesec}

\pgfplotsset{compat=1.18}

\definecolor{titleColor}{HTML}{800000}
\definecolor{blueColor}{HTML}{4682B4}
\definecolor{purpleColor}{HTML}{6A5ACD}
\definecolor{goldColor}{HTML}{DAA520}

\hypersetup{
    colorlinks=true,
    linkcolor=titleColor,
    urlcolor=blueColor
}

\pagestyle{fancy}
\fancyhf{}
\lhead{REEM V2}
\rhead{Demand Forecasting}
\cfoot{\thepage}

\titleformat{\section}
{\Large\bfseries\color{titleColor}}
{\thesection}{1em}{}

\titleformat{\subsection}
{\large\bfseries\color{blueColor}}
{\thesubsection}{1em}{}

\title{
\textcolor{titleColor}{
\textbf{REEM V2}\\[0.3cm]
Demand Forecasting and Model Validation
}
}

\author{
Pharmacy Supply Chain Intelligence\\
Research and Development Framework
}

\date{\today}

\begin{document}

\maketitle

\begin{center}
\begin{tikzpicture}[
    node distance=1.5cm,
    box/.style={
        rectangle,
        rounded corners,
        draw=titleColor,
        thick,
        minimum width=3.4cm,
        minimum height=0.9cm,
        align=center
    },
    arrow/.style={->,thick}
]

\node[box] (data) {Operational\\Synthetic Data};
\node[box,right=of data] (feature) {Demand\\Series};
\node[box,right=of feature] (model) {Forecast\\Models};
\node[box,below=of model] (validation) {Out-of-Sample\\Validation};
\node[box,left=of validation] (selection) {Model\\Selection};
\node[box,left=of selection] (management) {Management\\Interpretation};

\draw[arrow] (data) -- (feature);
\draw[arrow] (feature) -- (model);
\draw[arrow] (model) -- (validation);
\draw[arrow] (validation) -- (selection);
\draw[arrow] (selection) -- (management);

\end{tikzpicture}
\end{center}

\begin{abstract}
This report presents the first demand forecasting engine developed within
REEM V2. The engine is implemented entirely in pure Python without
third-party Python libraries. Four transparent forecasting methods are
evaluated for each medicine using an 80/20 temporal train--validation
framework. Model performance is assessed using Mean Absolute Error (MAE),
Root Mean Squared Error (RMSE), and forecast bias.

The underlying REEM V2 dataset is synthetic development data generated by
the IDEALBUILD environment. Therefore, the results presented here are for
model development and methodological validation only and must not be
interpreted as operational pharmacy recommendations.
\end{abstract}

\section{Purpose}

The purpose of this model is to establish a transparent baseline for
pharmacy demand forecasting.

The principal question is:

\begin{quote}
\textit{Given the observed demand history of a medicine, which simple
forecasting method provides the most accurate prediction of future demand?}
\end{quote}

The model is deliberately transparent. Rather than assuming that one
forecasting method is universally superior, REEM V2 evaluates several
methods and selects the best-performing method separately for each
medicine.

\section{Data Architecture}

The model uses the longitudinal synthetic pharmacy environment generated
by \texttt{idealbuild.py}.

\begin{center}
\begin{tikzpicture}[
    node distance=1.2cm,
    box/.style={
        rectangle,
        rounded corners,
        draw=blueColor,
        thick,
        minimum width=4.0cm,
        minimum height=0.8cm,
        align=center
    },
    arrow/.style={->,thick}
]

\node[box] (sales) {\texttt{schema/sales.csv}};
\node[box,below=of sales] (daily) {Daily Medicine\\Demand};
\node[box,below=of daily] (split) {80\% Training\\20\% Validation};
\node[box,below=of split] (forecast) {Four Forecasting\\Methods};
\node[box,below=of forecast] (metrics) {MAE / RMSE / Bias};
\node[box,below=of metrics] (select) {Best Model per\\Medicine};

\draw[arrow] (sales) -- (daily);
\draw[arrow] (daily) -- (split);
\draw[arrow] (split) -- (forecast);
\draw[arrow] (forecast) -- (metrics);
\draw[arrow] (metrics) -- (select);

\end{tikzpicture}
\end{center}

The current development environment contains:

\begin{itemize}[leftmargin=1.5cm]
\item 50 medicines
\item approximately 3.7 years of longitudinal observations
\item 47,383 synthetic sales transactions
\item four forecasting methods
\item 200 model evaluations
\item one selected forecasting model for each of the 50 medicines
\end{itemize}

\section{Forecasting Methods}

Four baseline methods are currently implemented.

\subsection{Naive Forecast}

The next observation is assumed to equal the most recent observation:

\[
\hat{D}_{t+1}=D_t
\]

This provides an important benchmark.

\subsection{Moving Average}

For a seven-day window:

\[
\hat{D}_{t+1}
=
\frac{1}{7}
\sum_{i=0}^{6}D_{t-i}
\]

The method smooths short-term fluctuations.

\subsection{Weighted Moving Average}

Recent observations receive greater weight:

\[
\hat{D}_{t+1}
=
\frac{
\sum_{i=1}^{k}iD_i
}{
\sum_{i=1}^{k}i
}
\]

This allows the model to respond more strongly to recent demand.

\subsection{Trend Forecast}

A simple least-squares linear trend is estimated over the recent
observation window and extrapolated one step forward.

The trend model is intentionally simple and interpretable.

\section{Validation Design}

The model uses a temporal holdout rather than randomly mixing observations.

\begin{center}
\begin{tikzpicture}[
    train/.style={
        rectangle,
        draw=blueColor,
        thick,
        minimum height=1cm,
        minimum width=7cm,
        align=center
    },
    test/.style={
        rectangle,
        draw=titleColor,
        thick,
        minimum height=1cm,
        minimum width=3cm,
        align=center
    }
]

\node[train] (train) {80\% Historical Data --- TRAINING};
\node[test,right=0cm of train] (test) {20\% --- TEST};

\draw[->,thick] (train.south) -- ++(0,-0.8)
node[midway,left]{Model estimation};

\draw[->,thick] (test.south) -- ++(0,-0.8)
node[midway,right]{Independent evaluation};

\end{tikzpicture}
\end{center}

This preserves the chronological structure of the demand series.

The validation period is never used to calculate the original training
forecast parameters.

\section{Evaluation Metrics}

\subsection{Mean Absolute Error}

\[
MAE=
\frac{1}{n}
\sum_{t=1}^{n}
|D_t-\hat{D}_t|
\]

MAE represents the average absolute forecasting error.

\subsection{Root Mean Squared Error}

\[
RMSE=
\sqrt{
\frac{1}{n}
\sum_{t=1}^{n}
(D_t-\hat{D}_t)^2
}
\]

RMSE gives greater influence to large forecasting errors.

\subsection{Forecast Bias}

\[
Bias=
\frac{1}{n}
\sum_{t=1}^{n}
(\hat{D}_t-D_t)
\]

Positive bias indicates systematic over-forecasting, while negative bias
indicates systematic under-forecasting.

\section{Current Results}

The model produced:

\begin{center}
\begin{tabular}{lr}
\toprule
Measure & Result\\
\midrule
Medicines evaluated & 50\\
Forecasting methods & 4\\
Total model evaluations & 200\\
Final medicine forecasts & 50\\
Validation design & 80/20 temporal holdout\\
Validation metrics & MAE, RMSE, Bias\\
Python external libraries & None\\
Data source & SIMULATED\_REEM\_V2\\
Data status & DEVELOPMENT\\
Operational use & PROHIBITED\\
\bottomrule
\end{tabular}
\end{center}

\section{Example: Acetylcysteine}

For medicine D0001, Acetylcysteine, the validation results were:

\begin{center}
\begin{tabular}{lrrr}
\toprule
Method & MAE & RMSE & Bias\\
\midrule
Naive & 3.8741 & 5.3597 & -0.0148\\
Moving Average & \textbf{2.9698} & \textbf{4.1041} & 0.0577\\
Weighted Moving Average & 2.9993 & 4.1690 & 0.0459\\
Trend & 3.2256 & 4.4701 & 0.0404\\
\bottomrule
\end{tabular}
\end{center}

The Moving Average model has the lowest MAE and RMSE among the four
methods for this medicine.

Therefore:

\[
\boxed{
\text{Selected model for D0001 = Moving Average}
}
\]

This illustrates an important REEM V2 principle:

\begin{center}
\textbf{Model selection is evidence-based and medicine-specific.}
\end{center}


\section{Analytical Results}

The validation dataset contains 200 model evaluations representing four
forecasting methods applied to 50 medicines. Model performance was assessed
using the chronological validation period rather than the training period.

\subsection{Forecast Accuracy by Method}

The mean validation MAE provides a useful overall comparison of the four
forecasting approaches. Lower values indicate smaller absolute forecasting
errors.

\begin{figure}[H]
\centering
\begin{tikzpicture}
\begin{axis}[
    ybar,
    bar width=18pt,
    width=0.92\textwidth,
    height=7cm,
    ylabel={Mean MAE},
    symbolic x coords={
        MOVING\_AVERAGE,
        WEIGHTED\_MOVING\_AVERAGE,
        TREND,
        NAIVE
    },
    xtick=data,
    xticklabel style={
        rotate=35,
        anchor=east,
        font=\small
    },
    ymin=0,
    ymax=1.15,
    nodes near coords,
    nodes near coords align={vertical},
    enlarge x limits=0.12,
    grid=major,
    major grid style={dashed},
]
\addplot table[
    x=Method,
    y=MeanMAE,
    col sep=comma
] {
Method,MeanMAE
MOVING\_AVERAGE,0.8221
WEIGHTED\_MOVING\_AVERAGE,0.8352
TREND,0.8826
NAIVE,1.0041
};
\end{axis}
\end{tikzpicture}
\caption{Mean validation MAE across forecasting methods. Lower values indicate better forecast accuracy.}
\label{fig:method-mae}
\end{figure}

The Moving Average method produced the lowest mean MAE at 0.8221,
followed closely by the Weighted Moving Average at 0.8352. The Trend method
recorded 0.8826, while the Naive method had the highest mean MAE at 1.0041.

This ranking indicates that, within the present synthetic development
environment, recent observed demand provides a stronger forecasting signal
than the simple persistence assumption or the fitted short-term trend.

\subsection{Model Selection Across Medicines}

The selected model is determined independently for each medicine using the
lowest validation MAE.

\begin{figure}[H]
\centering
\begin{tikzpicture}
\begin{axis}[
    ybar,
    bar width=24pt,
    width=0.88\textwidth,
    height=7cm,
    ylabel={Number of Medicines},
    symbolic x coords={
        MOVING\_AVERAGE,
        NAIVE,
        WEIGHTED\_MOVING\_AVERAGE,
        TREND
    },
    xtick=data,
    xticklabel style={
        rotate=35,
        anchor=east,
        font=\small
    },
    ymin=0,
    ymax=40,
    nodes near coords,
    enlarge x limits=0.12,
    grid=major,
    major grid style={dashed},
]
\addplot table[
    x=Model,
    y=Medicines,
    col sep=comma
] {
Model,Medicines
MOVING\_AVERAGE,35
NAIVE,9
WEIGHTED\_MOVING\_AVERAGE,4
TREND,2
};
\end{axis}
\end{tikzpicture}
\caption{Forecasting model selected for each of the 50 medicines.}
\label{fig:model-selection}
\end{figure}

The Moving Average model was selected for 35 of the 50 medicines,
representing 70\% of the development cohort. The Naive model was selected
for 9 medicines, the Weighted Moving Average for 4, and the Trend model for
2.

This result demonstrates an important REEM principle: model selection is
medicine-specific. The system does not assume that one forecasting method
must be optimal for every medicine.

\subsection{Highest Forecast-Demand Medicines}

The current selected forecasts identify medicines with the highest expected
daily demand under the development model.

\begin{figure}[H]
\centering
\begin{tikzpicture}
\begin{axis}[
    xbar,
    width=0.92\textwidth,
    height=9cm,
    xlabel={Selected Forecast (units/day)},
    symbolic y coords={
        D0001,
        D0018,
        D0027,
        D0043,
        D0021,
        D0042,
        D0048,
        D0016,
        D0005,
        D0009
    },
    ytick=data,
    xmin=0,
    xmax=10,
    nodes near coords,
    nodes near coords align={horizontal},
    enlarge y limits=0.06,
    grid=major,
    major grid style={dashed},
]
\addplot coordinates {
    (9.00,D0001)
    (9.00,D0018)
    (7.86,D0027)
    (7.57,D0043)
    (7.43,D0021)
    (7.43,D0042)
    (6.57,D0048)
    (5.86,D0016)
    (4.57,D0005)
    (4.57,D0009)
};
\end{axis}
\end{tikzpicture}
\caption{Top ten medicines by selected forecast demand.}
\label{fig:top-demand}
\end{figure}

The highest selected forecast is 9.00 units per day for D0001 and D0018.
The remaining medicines in the top ten range from 4.57 to 7.86 units per
day.

These values should be interpreted as model outputs from synthetic
development data. They are not procurement recommendations or estimates of
actual pharmacy demand.

\subsection{Analytical Interpretation}

Three findings emerge from the current experiment:

\begin{enumerate}
    \item The Moving Average method has the lowest overall mean validation
    MAE among the four candidate methods.

    \item Model performance is heterogeneous across medicines. The presence
    of four different selected models confirms that a single forecasting
    method is not uniformly optimal.

    \item Forecast magnitude identifies medicines requiring further
    analytical attention, but forecast magnitude alone is insufficient for
    inventory decisions. Future REEM stages must integrate forecast demand
    with lead time, stock position, safety stock, expiry, supplier
    reliability, and service-level requirements.
\end{enumerate}

The current result therefore represents a forecasting foundation rather than
a complete inventory optimisation system.

\section{Interpretation}

A forecasting model should not be judged only by the numerical forecast.
The quality of the forecasting process must also be considered.

For D0001, the Moving Average model achieved:

\[
MAE=2.9698
\]

This means that, on average, the absolute forecast error was approximately
2.97 units during the held-out validation period.

The RMSE was:

\[
RMSE=4.1041
\]

The difference between RMSE and MAE indicates that some larger errors
occurred during the validation period.

The bias was close to zero:

\[
Bias=0.0577
\]

Therefore, there is little evidence of strong systematic over- or
under-forecasting in this validation example.

\section{Model Selection Logic}

For every medicine, REEM V2 calculates the validation error for all four
methods.

The method with the smallest MAE becomes the selected baseline model.

\begin{center}
\begin{tikzpicture}[
    box/.style={
        rectangle,
        rounded corners,
        draw=purpleColor,
        thick,
        minimum width=4cm,
        minimum height=0.85cm,
        align=center
    },
    arrow/.style={->,thick}
]

\node[box] (m1) {Naive};
\node[box,below=0.5cm of m1] (m2) {Moving Average};
\node[box,below=0.5cm of m2] (m3) {Weighted Moving Average};
\node[box,below=0.5cm of m3] (m4) {Trend};

\node[box,right=4cm of m2] (compare) {Compare Validation\\MAE};

\node[box,right=4cm of compare] (best) {Select Lowest\\MAE};

\draw[arrow] (m1) -- (compare);
\draw[arrow] (m2) -- (compare);
\draw[arrow] (m3) -- (compare);
\draw[arrow] (m4) -- (compare);
\draw[arrow] (compare) -- (best);

\end{tikzpicture}
\end{center}

\section{Evidence Status}

The present results belong to the development stage.

\begin{center}
\begin{tikzpicture}[
    box/.style={
        rectangle,
        rounded corners,
        draw=goldColor,
        thick,
        minimum width=4cm,
        minimum height=0.8cm,
        align=center
    },
    arrow/.style={->,thick}
]

\node[box] (a) {Calculated};
\node[box,right=1.2cm of a] (b) {Analytically Useful};
\node[box,right=1.2cm of b] (c) {Validated};
\node[box,right=1.2cm of c] (d) {Operationally Trusted};

\draw[arrow] (a) -- (b);
\draw[arrow] (b) -- (c);
\draw[arrow] (c) -- (d);

\end{tikzpicture}
\end{center}

The current REEM V2 demand results have reached the calculation and
development-validation stages.

They have \textbf{not} reached operational trust because the underlying
dataset is synthetic.

\section{Limitations}

The current model has several limitations.

\begin{enumerate}[leftmargin=1.5cm]
\item The underlying dataset is simulated.
\item The model represents baseline forecasting rather than a complete
probabilistic forecasting system.
\item Intermittent demand requires dedicated methods.
\item Forecast uncertainty has not yet been fully modelled.
\item Real pharmacy data may contain stockout-induced demand censoring.
\item Calendar, promotion, substitution and external contextual variables
require further development.
\end{enumerate}

These limitations are intentional. REEM V2 is being developed
incrementally so that each modelling layer can be inspected and validated
before more complex methods are introduced.

\section{Next Development Stage}

The next modelling layers can build on the validated demand foundation:

\begin{center}
\begin{tikzpicture}[
    box/.style={
        rectangle,
        rounded corners,
        draw=blueColor,
        thick,
        minimum width=4.2cm,
        minimum height=0.75cm,
        align=center
    },
    arrow/.style={->,thick}
]

\node[box] (demand) {Validated Demand Engine};
\node[box,below=0.7cm of demand] (inter) {Intermittent Demand};
\node[box,below=0.7cm of inter] (rop) {Reorder Point};
\node[box,below=0.7cm of rop] (safe) {Safety Stock};
\node[box,below=0.7cm of safe] (eoq) {EOQ};
\node[box,below=0.7cm of eoq] (sim) {Monte Carlo};
\node[box,below=0.7cm of sim] (opt) {Inventory Optimisation};

\draw[arrow] (demand) -- (inter);
\draw[arrow] (inter) -- (rop);
\draw[arrow] (rop) -- (safe);
\draw[arrow] (safe) -- (eoq);
\draw[arrow] (eoq) -- (sim);
\draw[arrow] (sim) -- (opt);

\end{tikzpicture}
\end{center}

\section{Conclusion}

The first REEM V2 demand forecasting engine successfully demonstrates a
complete modelling cycle using pure Python:

\[
\text{Operational Structure}
\rightarrow
\text{Demand Series}
\rightarrow
\text{Forecast}
\rightarrow
\text{Validation}
\rightarrow
\text{Model Selection}
\rightarrow
\text{Interpretation}
\]

The framework evaluates multiple forecasting methods rather than assuming
that one model is universally appropriate.

The current results should be regarded as a controlled development
environment for building and validating the REEM V2 intelligence
architecture.

\begin{center}
\vspace{0.5cm}

\textcolor{titleColor}{
\Large\textbf{
Calculated $\rightarrow$ Validated $\rightarrow$ Trusted
}}

\vspace{0.3cm}

\textit{
Synthetic development evidence must not be confused with operational
pharmacy evidence.
}

\end{center}

\end{document}








\usepackage{booktabs}
\usepackage{soul}
\usepackage[most]{tcolorbox}
\usepackage{tikz}
\usetikzlibrary{arrows.meta, positioning, shapes}
\usepackage{sectsty}

